% !TEX root = main.tex
\section{Local Certainty and Duality}
\label{sec:duality}

The deterministic hull mixes states that are certain on all events at once. A
weaker requirement is certainty within each of a family of contexts. This section places the PCS between the deterministic hull
and the resulting local relaxation. It ends with a quantum event logic,
which has no certain states.

\subsection{The Local Relaxation}
\label{sec:duality-local}

\begin{definition}[local relaxation~\cite{acin2015}\leanref{local-calc}]
\label{def:local}
For a given finite set $E$ of events and a family $\mathcal{C}$ of
subsets of $E$, the \emph{contexts}, the \emph{local relaxation} is the set
of $x : E \to \mathbb{R}_{\ge 0}$ such that $\sum_{e \in C} x_e \le 1$
for every context $C$. The \emph{deterministic hull} is $\Mix$ of the
elements of the local relaxation with values in $\{0,1\}$. The
\emph{normalized} local relaxation and deterministic hull are the sets defined in the
same way with $=$ in place of $\le$.
\end{definition}

\noindent
The local relaxation contains its
$\{0,1\}$-valued elements and is closed under mixtures, so
Definition~\ref{def:gap} applies to it, with the deterministic hull as
$\Mix$ of its certain states. The normalized local relaxation and deterministic hull are
the general probabilistic models and the classical models of a family of
contexts~\cite{acin2015}.

For a formula $A$, take as contexts the supports of the certain states of
$A^{\perp}$. The pairing with a $\{0,1\}$-valued function is the sum over
its support, so the local relaxation of $A$ is the orthogonal of the certain
states of $A^{\perp}$\leanref{loc-linear}. We write $\loc(A)$ for it,
and we prove that it bounds the PCS from above.

\begin{theorem}[local certainty\leanref{loc-linear}\leanref{recipe-core}\leanref{recipe-formula}]
\label{thm:loc-linear}
For every formula $A$, the deterministic hull of $A$ is contained in $P(A)$, which is contained in $\loc(A)$. Moreover, $P(A) = \loc(A)$ iff $A^{\perp}$
has no gap.
At $T$ the second inclusion is an equality, and if $X, Y, X', Y'$ have at
least two elements each, the first is strict.
\end{theorem}

\noindent
An orthogonal does not change when a set is replaced by its mixtures, so
$\loc(A)$ is the orthogonal of the deterministic hull of $A^{\perp}$. By
Definition~\ref{def:mall}(iii), $P(A)$ is the orthogonal of $P(A^{\perp})$. The deterministic hull of $A^{\perp}$ is contained
in $P(A^{\perp})$, which gives $P(A) \subseteq \loc(A)$. The deterministic hull and
the PCS of every formula are equal to their biorthogonals, the
deterministic hull by the bipolar theorem in the proof of
Theorem~\ref{thm:second-exact}, which gives the equivalence. At $T$ the dual is the core, which has no gap
(Theorem~\ref{thm:second-exact}(i)), and the first inclusion is strict by
Theorem~\ref{thm:pentagon}.

\subsection{A Clique Outside the PCS}
\label{sec:duality-shannon}

The converse of the last statement of Theorem~\ref{thm:mall} fails: some
cliques of the coherence graph are not certain states. Let
$X = Y = X' = Y' = \mathrm{Bool}$. On the five points of
Theorem~\ref{thm:pentagon}, two distinct points cohere in $T$ iff they are
not exclusive. Order them as $c_0, \dots, c_4$ along the cycle of length
$5$ of the exclusivity graph. The coherence graph on them is the
complement of this cycle, again a cycle of length $5$. Two distinct points
of $T \otimes T$ are incoherent iff they are incoherent in one component,
so a clique of $(T \otimes T)^{\perp}$ contained in
$\{c_0, \dots, c_4\}^2$ is a stable set of the strong product of two
cycles of length $5$. The set $\{(c_i, c_{2i}) \mid i \in \mathbb{Z}/5\}$
is such a stable set, of size $5$, Shannon's stable
set~\cite{lovasz1979}. We show that it gives a clique outside the
PCS.

\begin{theorem}[a clique outside the PCS\leanref{shannon-clique}\leanref{shannon-not-mem}\leanref{tensorgap-bool}]
\label{thm:shannon}
Shannon's set is a clique of the coherence graph of $(T \otimes T)^{\perp}$,
and its indicator does not lie in $P((T \otimes T)^{\perp})$.
\end{theorem}

\noindent
The function $z$ of Theorem~\ref{thm:pentagon}, with value $\tfrac12$ on
the five points, lies in $P(T)$, so $z \otimes z$ lies in $P(T \otimes T)$.
Its pairing with the indicator of Shannon's set is
$5 \cdot \tfrac14 = \tfrac54 > 1$.

The indicator of every clique of $A$ lies in $\loc(A)$. The certain
states of $A^{\perp}$ are cliques of $A^{\perp}$ (Theorem~\ref{thm:mall}). A clique of $A$ and a
clique of $A^{\perp}$ share at most one point,
since two distinct points cannot cohere in both. This allows us to prove
that Shannon's set also separates the PCS from the local
relaxation.

\begin{corollary}[the second inclusion\leanref{loc-linear}]
\label{cor:tt}
At $(T \otimes T)^{\perp}$, the PCS is strictly contained in
the local relaxation.
\end{corollary}

\noindent
The first inclusion is strict there too, by Theorem~\ref{thm:c4}: the
points $((t,t),(t,t))$, $((t,t),(t,f))$, $((f,t),(t,t))$, $((f,t),(t,f))$
form an induced $4$-cycle of $T$\leanref{tt-gap}.
We did not find the placement of the PCS between the
deterministic hull and the local relaxation
(Theorem~\ref{thm:loc-linear}) in the literature.

\subsection{Negation, Tightness, and the Tensor}
\label{sec:duality-tight}

Two sets $U$ and $V$ of weight vectors on one web form a \emph{tight
pair} if each is the orthogonal of the other, $U = V^{\perp}$ and
$V = U^{\perp}$~\cite{hyland-schalk2003}. A PCS and its
orthogonal form a tight pair. The question is when the deterministic hulls of a formula and of its negation form a tight
pair, and we prove the following answer, whose part (i) extends the
antiblocking duality below.

\begin{theorem}[duality\leanref{orth-recipe}\leanref{tight}]
\label{thm:duality}
Let $A$ be a formula.
(i) The orthogonal of the deterministic hull of $A$ is $\loc(A^{\perp})$, and the
orthogonal of $\loc(A)$ is the deterministic hull of $A^{\perp}$.
(ii) The deterministic hulls of $A$ and of $A^{\perp}$ form a tight pair iff the
deterministic hull of $A$ is $\loc(A)$, iff neither $A$ nor $A^{\perp}$ has a gap.
(iii) If $X, Y, X', Y'$ have at least two elements each, the deterministic hulls of
the core $K$ and of $T$ do not form a tight pair.
\end{theorem}

\noindent
Part (i) combines the biorthogonal closure of the deterministic hull with the fact
that an orthogonal does not change when a set is replaced by its
mixtures. Each of the two equations of a tight pair of deterministic hulls implies the
other, since the deterministic hull is equal to its biorthogonal. Part (ii) follows
with Theorem~\ref{thm:loc-linear}, and (iii)
with Theorem~\ref{thm:second-exact}(i) and the gap at $T$
(Theorem~\ref{thm:pentagon}).

When the certain states of $A$ and of $A^{\perp}$ are the cliques of their
coherence graphs, (i) is the antiblocking identity between the stable set
polytope of a graph and the clique-constrained polytope of its
complement, which joins two distinct vertices iff the graph does
not~\cite{fulkerson1971,knuth1994}. In that case tightness holds iff the
coherence graph is perfect~\cite{lovasz1972,chvatal1975}. For a coherence space $E$ in the sense of Girard, let $p(E)$ be the set of $x \ge 0$ with total
weight at most $1$ on every clique of $E^{\perp}$. Then
$p(E^{\perp}) = p(E)^{\perp}$ when $E$ is generated from $\mathbf{1}$ by
$\perp$ and $\mathbin{\&}$, so that its coherence graph is a cograph, and
not when its coherence graph has an induced cycle of length
$5$~\cite{ehrhard-td2}. A family of finite sets and its orthogonal family
have dual norm balls iff they are the cliques and the stable sets of a
perfect graph~\cite{bfp2026}. These results concern the cliques of a coherence graph.
Theorem~\ref{thm:duality} concerns its certain states, which can be fewer
than the cliques (Theorem~\ref{thm:shannon}). To our knowledge, part (ii), the equivalence of tightness with the
absence of a gap at $A$ and at $A^{\perp}$, does not appear in the
literature.

We show that tightness holds at first order and fails at a tensor of
first-order types.

\begin{corollary}[tightness and the tensor\leanref{tight-tensor}]
\label{cor:tensor}
For all $n$ and $m$, the local relaxation of $\mathbf{n} \multimap
\mathbf{m}$ is $P(\mathbf{n} \multimap \mathbf{m})$, so the deterministic hulls of
$\mathbf{n} \multimap \mathbf{m}$ and of its negation form a tight pair.
For $n, m, n', m' \ge 2$, the deterministic hulls of $(\mathbf{n} \multimap
\mathbf{m}) \otimes (\mathbf{n'} \multimap \mathbf{m'})$ and of its
negation do not.
\end{corollary}

\noindent
The indicator of a row $\{a\} \times \{1, \dots, m\}$ is a certain state
of the negation of $\mathbf{n} \multimap \mathbf{m}$. These indicators cut
out the row-substochastic matrices, which form
$P(\mathbf{n} \multimap \mathbf{m})$. With
Proposition~\ref{prop:first-order} the deterministic hull is the local relaxation. The
second statement is Theorem~\ref{thm:duality} with the gap of
Theorem~\ref{thm:pentagon}. So tightness at $A$ and at $B$ does not
imply tightness at $A \otimes B$, although $P(A \otimes B)$ is determined
by $P(A)$ and $P(B)$ (Definition~\ref{def:mall}(iii)).

\subsection{Bell's Theorem at the Core}
\label{sec:duality-bell}

Let $X = Y = X' = Y' = \mathbf{2}$. A point $((x, a), (y, b))$ of the web
of the core $K$ is the event that two parties with inputs $x$ and $y$ give
outputs $a$ and $b$. Two events are \emph{locally orthogonal} if some
party has equal inputs and different outputs. For distinct events this is the negation of exclusivity
(Section~\ref{sec:second-pentagon}). An element of $P(K)$ is \emph{normalized} if it has total
weight $1$ for each pair of inputs. The second statement of the following
form of Bell's theorem is known~\cite{bell1964,chsh1969,popescu-rohrlich1994},
the first restates Theorem~\ref{thm:qstab}, and we formalize both.

\begin{theorem}[Bell at the core~\cite{bell1964,popescu-rohrlich1994}\leanref{bell}]
\label{thm:bell}
Every certain state of $T$ is the indicator of a set of pairwise locally
orthogonal events. The Popescu--Rohrlich box~\cite{popescu-rohrlich1994},
with value $\tfrac12$ on the
events with $a + b \equiv xy \pmod 2$ and $0$ elsewhere, lies in
$\loc(K)$ and not in $P(K)$.
\end{theorem}

\noindent
Two events that are not locally orthogonal are both produced by one
deterministic product strategy, a pair of deterministic programs, one for
each party, which lies in $P(K)$ and has pairing $2$ with their
indicator. No three events in the support of the box are pairwise
locally orthogonal, so the box has total weight at most $1$ on the support
of every certain state of $T$, and lies in $\loc(K)$. By
Theorem~\ref{thm:second-exact}, $P(K)$ is the deterministic hull of $K$. Every
certain state of $K$ lies below a deterministic product strategy, which
has total weight at most $3$ on the eight events with $a + b \equiv xy$,
the CHSH inequality~\cite{chsh1969}. The box has total weight $4$ there.
By Theorem~\ref{thm:loc-linear}, $P(K) \ne \loc(K)$ iff $T$ has a gap.

The normalized elements of $P(K)$ are the Bell-local behaviours, i.e. the
mixtures of deterministic product strategies that always terminate. A set
of pairwise locally orthogonal events meets each clique of $K$ in at most
one event, so $\loc(K)$ is the set of functions with total weight at most
$1$ on every such set. Its normalized elements are the no-signalling
behaviours~\cite{fritz-lo2013}. These two identifications are not
formalized.
